\section{沿一条 Shell 命令回看整个系统}

在 Shell 中输入 \code{cat /tmp/note}，表面上只是读取并打印一个文件。这个
命令恰好能把前十九章重新连在一起：键盘产生中断，TTY 提交一行输入，Shell
解析参数并启动外部 ELF，用户程序通过系统调用打开路径，VFS 穿过
\filepath{/tmp} 挂载点，memfs 返回字节，最后标准输出又回到终端。

\begin{center}
\begin{tikzpicture}[os diagram]
  \node[os state] (key) {键盘 IRQ1};
  \node[os block,right=of key] (tty) {TTY\\提交一行};
  \node[os block,right=of tty] (shell) {Shell\\解析与启动};
  \node[os state,right=of shell] (cat) {cat ELF\\Ring 3};
  \node[os block,below=of cat] (syscall) {open/read/write\\系统调用};
  \node[os block,left=of syscall] (vfs) {VFS\\路径与挂载};
  \node[os state,left=of vfs] (memfs) {memfs\\文件字节};
  \draw[os arrow] (key) -- (tty);
  \draw[os arrow] (tty) -- (shell);
  \draw[os arrow] (shell) -- (cat);
  \draw[os arrow] (cat) -- (syscall);
  \draw[os arrow] (syscall) -- (vfs);
  \draw[os arrow] (vfs) -- (memfs);
\end{tikzpicture}
\end{center}

\topic{命令执行前，启动链已经做了什么}
CPU 从 ROM 复位向量开始，固件建立串口并读取 Stage 1。Stage 1 打开 A20，
进入保护模式和长模式，检查并装入 Kernel ELF。Kernel 接管 GDT、IDT、TSS、
页表和中断控制器，挂载文件系统，再由 PID 1 从磁盘启动 Shell。Shell 不是
内核函数，而是一个普通用户进程。

\topic{输入一行时，异步事件怎样进入普通代码}
键盘控制器在端口准备好扫描码后触发 IRQ1。中断入口保存现场，驱动读取字节，
TTY 把 make/break 和控制字符解释为行编辑动作。回车使等待输入的 Thread
变为 Ready；调度器随后让 Shell 从 read 系统调用继续。

\topic{启动 cat 时，进程与文件怎样组合}
Shell 准备 argv 和文件描述符，创建子进程并执行磁盘上的 cat ELF。装载器先
建立 VMA；代码页第一次取指时再从文件后备读入。标准输出 fd=1 继续引用终端
打开实例，因此 cat 不需要知道串口或显存地址。

\topic{读取路径时，各层分别拥有哪件事}
\begin{osgridlongtable}{L{0.21\textwidth}L{0.31\textwidth}L{0.36\textwidth}}
\textbf{层} & \textbf{在本例中保存什么} & \textbf{不处理什么}\\ \hline
FileTable & 当前进程的 fd 编号与 fd flags & 路径分量和磁盘布局\\ \hline
FileDescription & offset、读写状态与共享引用 & 当前工作目录\\ \hline
VFS & \filepath{/tmp/note} 的挂载切换与 vnode & memfs 字节怎样存放\\ \hline
memfs & 节点、名称和文件内容 & ATA 端口与 journal\\ \hline
TTY & 输出队列与终端归属 & 文件路径\\ \hline
\end{osgridlongtable}

\topic{换成持久文件以后，最后一段路径发生变化}
若读取 \filepath{/home/note}，VFS 仍完成相同的分量解析，但后端变为 rootfs。
数据可能先从文件页缓存命中；缺失时形成块请求并等待 IRQ14。写入后还要经过
writeback、journal 顺序和 ATA flush，才能讨论断电后的结果。用户程序的
read/write 接口没有因此改变。

\topic{用一条命令检查自己的理解}
读完本书后，可以画出这条命令涉及的 Thread 状态、CPL、CR3、fd 引用、VMA、
Vnode 和中断向量。再任选一个失败点，例如 cat 的 ELF 被截断或
\filepath{/tmp/note} 不存在，说明错误在哪一层产生、沿哪条返回路径到达
Shell。能把成功与失败都走通，比记住模块名单更接近真正理解系统。
